% CSC5515 G14 final report -- DRAFT.
% Move the body into the course LaTeX template in Week 12. Every "\todo{...}" needs a GPU result or a
% decision; everything else is taken from the proposal or computed from the real data (docs/week7_data_report.md).
\documentclass[11pt]{article}
\usepackage[margin=2.5cm]{geometry}
\usepackage{fontspec}
\usepackage{polyglossia}
\setdefaultlanguage{english}
\setotherlanguage{bengali}
\newfontfamily\bengalifont[Script=Bengali]{Noto Sans Bengali}
\usepackage{booktabs}
\usepackage{xcolor}
\usepackage{hyperref}
\newcommand{\todo}[1]{\textcolor{red}{[TODO: #1]}}
\newcommand{\bn}[1]{\textbengali{#1}}

\title{Fine-Tuning an LLM to Translate Bangla Regional Dialects into Standard Bangla}
\author{Esm-e Moula Chowdhury Abha (26-94089-2)\\CSC5515 Natural Language Processing, Group G14}
\date{December 2026}

\begin{document}
\maketitle

\begin{abstract}
\todo{Write last. One sentence each: task, method (QLoRA on Qwen2.5-7B-Instruct), data (Vashantor, ONUBAD),
main result (chrF and human rating vs.\ prompting), dialect finding (RQ2), transfer finding (RQ3).}
\end{abstract}

\section{Introduction}
Bangla is spoken in many regional dialects. Chittagonian and Sylheti differ so much from standard Bangla
that many speakers from other regions cannot follow them: the same sentence can use different words, verb
endings and pronouns while its meaning stays the same. Large language models are trained mostly on standard
written Bangla and often misread dialect text in social-media posts, customer messages or call-centre
transcripts. A model that rewrites dialect text into standard Bangla without changing its meaning lets
existing Bangla tools work on such text. This project fine-tunes a small open-weight LLM for this task and
measures how well meaning is preserved.

\begin{itemize}
\item \textbf{RQ1.} Does a QLoRA-tuned LLM preserve meaning better than zero-shot and few-shot prompting of
the same model, measured by chrF, LaBSE similarity and human ratings?
\item \textbf{RQ2.} Which dialects are hardest, and is the difficulty explained by word-level distance from
standard Bangla, defined as 100 minus the chrF of the copy baseline?
\item \textbf{RQ3.} Can a model trained on four dialects translate a fifth dialect it has not seen?
\end{itemize}

\section{Related work}
Faria et al.\ \cite{vashantor} released Vashantor, a parallel corpus of five dialects and standard Bangla,
and fine-tuned mT5 and BanglaT5 \cite{banglat5}, reporting BLEU from 36.75 (Chittagong) to 69.06
(Mymensingh). Jawad et al.\ \cite{jawad} benchmarked zero- and few-shot prompting of LLMs on 600 YouTube
comments in four dialects and found that transliteration helps closed models. Sami et al.\ \cite{sami}
used retrieval-augmented prompting for the opposite direction, standard-to-dialect. Sultana et al.\
\cite{onubad} collected ONUBAD, a separate corpus of Chittagong, Sylhet and Barishal sentences, which we use
as an out-of-domain test set. We fine-tune instruction-tuned LLMs with QLoRA \cite{lora,qlora}, measure
meaning with chrF \cite{chrf}, LaBSE \cite{labse} and human ratings, and add leave-one-dialect-out transfer.

\section{Data}
\paragraph{Vashantor.} Mendeley Data V2 (doi:10.17632/bj5jgk878b.2, CC BY 4.0), CSV release. Each file pairs
a dialect sentence (\texttt{<dialect>\_bangla\_speech}) with its standard sentence (\texttt{bangla\_speech});
Banglish, English and region columns are not used. The official split has 1,875 / 250 / 375 pairs per
dialect. Cleaning (Unicode NFC, zero-width characters removed, whitespace collapsed) removed no pairs:
the release has no empty and no exact duplicate pairs.

\paragraph{Alignment of splits.} The same standard sentences are translated into all five dialects, so a
test sentence must not appear in training through another dialect. Pairs are grouped by their standard
sentence, ignoring punctuation and spaces, and a group found in more than one split is moved entirely to its
test or validation split. One standard sentence (\bn{আপনি কেমন আছেন?}) occurred in both the training and
the validation split; its 5 training pairs were moved, giving 1,874 / 251 / 375 pairs per dialect
(Table~\ref{tab:data}). Seventeen groups have fewer than five members because the standard sentence is spelled
differently in some dialect files (e.g.\ \bn{নিষেধ}/\bn{নিশেধ}); each lies in a single split. Five validation
and five test standard sentences differ from a training sentence by one word (e.g.\ \bn{আমার}/\bn{তোমার
বন্ধুরা}~\ldots); they are different sentences and are kept.

\begin{table}[h]\centering
\begin{tabular}{lrrrr}\toprule
Split & Pairs & Per dialect & Standard sentences & Mean words \\\midrule
Train & 9,370 & 1,874 & 1,876 & 6.7 \\
Validation & 1,255 & 251 & 250 & 5.5 \\
Test & 1,875 & 375 & 381 & 7.3 \\\bottomrule
\end{tabular}
\caption{Vashantor after cleaning and alignment (mean words of the standard sentence).}\label{tab:data}
\end{table}

\paragraph{ONUBAD.} Mendeley Data V2 (doi:10.17632/6ft99kf89b.2, CC BY 4.0). Only the sentence file is used
(the word and clause files are not). After removing 144 duplicate pairs there are 2,796 pairs (Chittagong
935, Sylhet 931, Barishal 930). \textbf{323 pairs have a standard sentence that also occurs in the Vashantor
training set} and 126 a dialect sentence that does, so they are not really out-of-domain. We therefore report
ONUBAD both in full and as the 2,437 \emph{novel} pairs without this overlap.

\section{Method}
\paragraph{Models.} Main model: Qwen2.5-7B-Instruct \cite{qwen}; second model: Llama-3.1-8B-Instruct
\cite{llama}. Llama-3.3-70B-Instruct is run zero-shot through the Groq API as a reference for a much larger
untuned model.

\paragraph{Baselines.} (i) \emph{Copy}: the input is output unchanged. (ii) \emph{Dictionary}: per dialect,
word replacements are learned from training pairs aligned with a longest-matching-block alignment of the two
token sequences; each dialect word is replaced by its most frequent standard form, unknown words are copied.
(iii) Zero-shot prompting with the instruction below (three paraphrases = three runs) and (iv) 5-shot prompting
with examples from the same dialect's training split (three draws).

\begin{quote}\small\ttfamily
Rewrite the following sentence from the \{dialect\} dialect into standard Bangla.\\
Keep the meaning unchanged. Output only the standard Bangla sentence.\\
Dialect sentence: \{text\}\\
Standard Bangla:
\end{quote}

\paragraph{QLoRA fine-tuning.} 4-bit NF4 base model, LoRA rank 32, $\alpha=64$, dropout 0.05 on all linear
layers, learning rate $2\cdot10^{-4}$ with cosine schedule and 3\% warm-up, effective batch 16, 2 epochs,
maximum length 256 tokens, loss on the answer tokens only; the same chat template is used for prompting and
fine-tuning. Validation chrF on a fixed stratified 250-pair validation subset is measured twice per epoch and
the best adapter is kept. Llama-3.1-8B uses the same settings. BanglaT5 is fully fine-tuned (fp32, learning rate
$3\cdot10^{-4}$, batch 16, best of 8 epochs by validation chrF) on inputs of the form
``\texttt{<Dialect>: <sentence>}'' normalised with the BanglaT5 normaliser.

\paragraph{Decoding.} Greedy, at most 128 new tokens, stopping at the end-of-turn token. Empty or unusable
outputs are kept and scored 0.

\paragraph{Runs.} \todo{Final run table with actual GPU hours from checkpoints/*/log.json.}

\section{Evaluation}
chrF (sacreBLEU) is the primary metric; BLEU is reported for comparison with \cite{vashantor}; LaBSE cosine
similarity between output and reference is the secondary meaning measure. Results are mean $\pm$ std over
three runs, with 95\% paired bootstrap confidence intervals (1,000 resamples) for differences. Dialect
distance is $100-\mathrm{chrF}$ of the copy baseline, with the copy baseline's word error rate as a check.
150 test outputs are rated 1--5 for meaning, blind to the system; a second annotator rates 50 of them and
agreement is quadratic-weighted Cohen's $\kappa$.

\section{Results}
\subsection{Baselines and dialect distance (validation)}
\begin{table}[h]\centering\small
\begin{tabular}{lrrrrrrrr}\toprule
System & Chit. & Noak. & Sylh. & Bari. & Myme. & chrF & BLEU & LaBSE\\\midrule
Copy & 26.00 & 37.56 & 39.79 & 42.61 & 58.67 & 40.95 & 10.43 & 0.67\\
Dictionary & 59.66 & 65.82 & 65.49 & 76.01 & 87.96 & 70.97 & 44.78 & 0.91\\\bottomrule
\end{tabular}
\caption{Non-LLM baselines on Vashantor validation (chrF per dialect and overall).}\label{tab:baselines}
\end{table}

Dialect distance on validation: Chittagong 74.0 (copy WER 90.8), Noakhali 62.4 (72.7), Sylhet 60.2 (73.9),
Barishal 57.4 (65.1), Mymensingh 41.3 (45.7). The ranking matches the earlier finding that Chittagonian is the
hardest dialect \cite{vashantor}. The dictionary baseline is already strong (71.0 chrF), so the LLM systems must
be compared against it, not only against the copy baseline.

\subsection{Preliminary CPU runs (validation)}
Before the GPU runs, two smaller systems were run on CPU to test the pipeline end to end
(Table~\ref{tab:cpu}). BanglaT5 after a single epoch (63 minutes on 4 CPU cores) is already close to the
dictionary baseline but 1.1 chrF below it (95\% CI $[-2.2, -0.1]$). Qwen2.5-3B-Instruct without fine-tuning
does not beat copying the input: zero-shot is 1.1 chrF below the copy baseline (CI $[-1.6, -0.6]$), and
5-shot examples add 1.1 chrF over zero-shot (CI $[+0.6, +1.7]$), which only brings it level with copying (+0.1, CI $[-0.6, +0.7]$). Its outputs are rarely copies (1--7\%) and almost never English,
but they keep many dialect words, change pronouns or tense, or invent non-words. Even so, the copy-based
distance predicts its zero-shot per-sentence chrF closely (Spearman $\rho=-0.90$, CI $[-0.91, -0.89]$). These runs
motivate fine-tuning (RQ1) and show that prompting a small model is not enough.

\begin{table}[h]\centering\small
\begin{tabular}{lrrrrrrrr}\toprule
System & Chit. & Noak. & Sylh. & Bari. & Myme. & chrF & BLEU & LaBSE\\\midrule
BanglaT5, 1 epoch & 55.09 & 64.71 & 67.34 & 75.53 & 86.41 & 69.86 & 46.65 & 0.91\\
Qwen2.5-3B zero-shot (3 runs) & 27.98 & 35.28 & 40.54 & 40.46 & 54.97 & 39.87$\pm$0.76 & 9.52 & 0.68\\
Qwen2.5-3B 5-shot (3 runs) & 29.63 & 36.84 & 41.38 & 41.53 & 55.56 & 41.01$\pm$0.35 & 10.09 & 0.70\\\bottomrule
\end{tabular}
\caption{Preliminary CPU runs on Vashantor validation. Not final systems.}\label{tab:cpu}
\end{table}

\subsection{Main comparison (RQ1)} \todo{results/test/main\_table.md and significance.md.}
\subsection{Which dialects are hardest (RQ2)} \todo{results/test/rq2.md: Spearman $\rho$ with CI.}
\subsection{Leave-one-dialect-out (RQ3)} \todo{results/lodo/test/rq3\_lodo.md.}
\subsection{Ablations} \todo{results/ablations/test/main\_table.md.}
\subsection{Out-of-domain and contamination}
Baselines on ONUBAD (all 2,796 pairs / 2,437 novel pairs): copy 37.3 / 37.6 chrF, dictionary 60.8 / 59.5
chrF (LaBSE 0.86 / 0.85). The dictionary loses about 10 chrF out of domain, because many ONUBAD words do not
occur in Vashantor training. \todo{results/onubad*/main\_table.md; compare zero-shot on
Vashantor and on novel ONUBAD.}
\subsection{Human evaluation} \todo{human\_eval/human\_eval\_results.md.}
\subsection{Error analysis} \todo{error\_analysis/error\_summary.md with examples.}

\section{Discussion and limitations}
\paragraph{Contamination.} Vashantor has been public since 2023, so the pretrained models may have seen its
test sentences. ONUBAD (2025) is used as a check, but it is not fully independent of Vashantor: 323 of its
2,796 pairs share a standard sentence with Vashantor training data and 126 share a dialect sentence. All
ONUBAD results are therefore also reported on the 2,437 novel pairs. If a system's advantage shrinks much more
on novel ONUBAD than the baselines' does, memorisation is a likely cause.

\paragraph{Data.} Vashantor's standard sentences are shared across the five dialects, so the split is made by
standard sentence; this prevents the same meaning from appearing in training and test, but it also means the
test set covers a limited set of topics and short conversational sentences. ONUBAD covers only three dialects.
Some dialect sentences differ from standard Bangla only in spelling, and a few validation and test pairs are
near-copies, which favour conservative systems.

\paragraph{Metrics and annotation.} chrF rewards surface overlap and can miss meaning errors such as changed
negation or tense; LaBSE and the human ratings are meant to catch these. Human ratings come from one annotator,
with a second annotator on only 50 items, so the agreement estimate has a wide confidence interval.

\paragraph{Runs.} The main systems use three runs, but the ablations and the leave-one-dialect-out folds use a
single seed because of the GPU budget, so small differences there should not be over-interpreted. Llama-3.3-70B
is evaluated on a 500-pair subset only.

\begin{thebibliography}{11}\small
\bibitem{vashantor} F.~T.~J. Faria, M.~B. Moin, A.~A. Wase, M.~Ahmmed, M.~R. Sani, T.~Muhammad. Vashantor: A
large-scale multilingual benchmark dataset for automated translation of Bangla regional dialects to Bangla
language. arXiv:2311.11142, 2023.
\bibitem{banglat5} A.~Bhattacharjee et al. BanglaNLG and BanglaT5: Benchmarks and resources for evaluating
low-resource natural language generation in Bangla. Findings of EACL, 2023.
\bibitem{lora} E.~J. Hu et al. LoRA: Low-rank adaptation of large language models. ICLR, 2022.
\bibitem{qlora} T.~Dettmers, A.~Pagnoni, A.~Holtzman, L.~Zettlemoyer. QLoRA: Efficient finetuning of quantized
LLMs. NeurIPS, 2023.
\bibitem{chrf} M.~Popović. chrF: character n-gram F-score for automatic MT evaluation. WMT, 2015.
\bibitem{labse} F.~Feng, Y.~Yang, D.~Cer, N.~Arivazhagan, W.~Wang. Language-agnostic BERT sentence embedding.
ACL, 2022.
\bibitem{onubad} N.~Sultana, R.~Yasmin, B.~Mallik, M.~S. Uddin. ONUBAD: A comprehensive dataset for automated
conversion of Bangla regional dialects into standard Bengali dialect. Data in Brief, 58:111276, 2025.
\bibitem{jawad} M.~M. Jawad et al. Benchmarking large language models on Bangla dialect translation and
dialectal sentiment analysis. BLP-2025, pp.~322--337, 2025.
\bibitem{sami} K.~M.~J. Sami, D.~Sumit, A.~Hossain, F.~Sadeque. A comparative analysis of retrieval-augmented
generation techniques for Bengali standard-to-dialect machine translation using LLMs. arXiv:2512.14179, 2025.
\bibitem{qwen} Qwen Team. Qwen2.5 technical report. arXiv:2412.15115, 2024.
\bibitem{llama} A.~Grattafiori et al. The Llama 3 herd of models. arXiv:2407.21783, 2024.
\end{thebibliography}
\end{document}
